\section*{Chapitre \ref{ch:g10:storms-and-cyclones} --- Orages et cyclones tropicaux}

\begin{solution}{exo:g10:storms-and-cyclones:1}
Un \omterm{def:g10:storms-and-cyclones:thunderstorm}{orage} est une perturbation produite par un cumulonimbus, avec \omterm{def:g10:storms-and-cyclones:lightning}{éclairs} et
tonnerre, en général de fortes averses et des rafales, parfois de la grêle.
Un \omterm{def:g10:storms-and-cyclones:lightning}{éclair} est une gigantesque étincelle électrique qui décharge des régions
chargées de l'\omterm{def:g10:storms-and-cyclones:thunderstorm}{orage}. Le tonnerre est le bruit de l'air du canal de l'\omterm{def:g10:storms-and-cyclones:lightning}{éclair},
chauffé si brusquement qu'il se dilate avec un choc.
\end{solution}

\begin{solution}{exo:g10:storms-and-cyclones:2}
Un \omterm{def:g10:storms-and-cyclones:tropical-cyclone}{cyclone tropical} est une tempête large de centaines de kilomètres, formée
au-dessus d'une mer tropicale chaude autour d'un centre de basse pression,
avec des vents qui spiralent vers l'intérieur et des bandes orageuses. L'\omterm{def:g10:storms-and-cyclones:tropical-cyclone}{œil}
est son centre calme, presque sans nuages~; le \omterm{def:g10:storms-and-cyclones:tropical-cyclone}{mur de l'œil} est l'anneau
d'\omterm{def:g10:storms-and-cyclones:thunderstorm}{orages} imposants qui entoure l'\omterm{def:g10:storms-and-cyclones:tropical-cyclone}{œil}. Les vents sont les plus forts dans le
\omterm{def:g10:storms-and-cyclones:tropical-cyclone}{mur de l'œil}.
\end{solution}

\begin{solution}{exo:g10:storms-and-cyclones:3}
\num{70}~nœuds~: catégorie 1. \num{120}~nœuds~: catégorie 4.
\num{140}~nœuds~: catégorie 5. Les deux derniers sont des ouragans majeurs
(catégorie 3 ou plus).
\end{solution}

\begin{solution}{exo:g10:storms-and-cyclones:4}
$343 \times 9 \approx \qty{3.1}{km}$, puis $343 \times 15 \approx
\qty{5.1}{km}$. Le second \omterm{def:g10:storms-and-cyclones:lightning}{éclair} est plus loin~: l'\omterm{def:g10:storms-and-cyclones:thunderstorm}{orage} s'éloigne.
\end{solution}

\begin{solution}{exo:g10:storms-and-cyclones:5}
La foudre tend à frapper le point le plus haut des environs~: un arbre isolé
dans un champ est ce point, et le courant de la décharge le parcourt puis se
répand dans le sol tout autour. Il faut entrer dans un bâtiment ou dans une
voiture fermée, et y rester bien après le dernier coup de tonnerre.
\end{solution}

\begin{solution}{exo:g10:storms-and-cyclones:6}
Un nœud vaut $1852/3600 \approx \qty{0.514}{m/s}$. \num{64}~nœuds~: $64 \times
0.514 \approx \qty{32.9}{m/s}$, et $32.9 \times 3.6 \approx
\qty{119}{km/h}$. \num{137}~nœuds~: $\qty{70.5}{m/s}$, environ
\qty{254}{km/h} (le tableau officiel, arrondi en passant par une autre
unité, donne \qty{252}{km/h}).
\end{solution}

\begin{solution}{exo:g10:storms-and-cyclones:7}
$44 \times 86\,400 \approx \num{3.8e6}$ \omterm{def:g10:storms-and-cyclones:lightning}{éclairs} par jour, et $3.8\times10^6
\times 365.25 \approx \num{1.4e9}$, environ 1.4~milliard par an.
\end{solution}

\begin{solution}{exo:g10:storms-and-cyclones:8}
$\Delta p = 1010 - 935 = \qty{75}{hPa}$~: un bombement d'environ \qty{75}{cm}.
\end{solution}

\begin{solution}{exo:g10:storms-and-cyclones:9}
Ils se forment au-dessus des océans tropicaux chauds, surtout entre
\qty{10}{\degree} et \qty{30}{\degree} de latitude dans les deux hémisphères,
avec très peu de cas dans l'Atlantique Sud. Ils dérivent vers l'ouest, puis
s'incurvent vers les pôles, où ils s'affaiblissent au-dessus d'eaux plus
froides. Ils sont absents de l'équateur parce que la rotation de la Terre
n'y imprime aucune déviation à l'air qui afflue~: sans elle, l'air file
droit dans la basse pression et la comble, et aucune tempête tournante ne
peut s'amorcer.
\end{solution}

\begin{solution}{exo:g10:storms-and-cyclones:10}
$1031 - 1.56 \times 70 \approx \qty{922}{hPa}$~; $1031 - 1.56 \times 40
\approx \qty{969}{hPa}$.
\end{solution}

\begin{solution}{exo:g10:storms-and-cyclones:11}
Sur la terre ferme, la tempête perd son apport d'eau qui s'évapore, et avec
lui la \omterm{def:g10:earth-energy-budget:latent-heat}{chaleur latente} qui l'entretient~; la surface rugueuse freine aussi
les vents. Au-dessus d'une eau froide, l'évaporation est bien plus faible,
car l'air froid contient bien moins de vapeur d'eau à saturation~: l'apport
de \omterm{def:g10:earth-energy-budget:latent-heat}{chaleur latente} s'amenuise.
\end{solution}

\begin{solution}{exo:g10:storms-and-cyclones:12}
À \qty{0}{\degreeCelsius}~: $\sqrt{1.4 \times 8.314 \times 273.15/0.02897}
\approx \qty{331}{m/s}$~; à \qty{30}{\degreeCelsius}, avec $T =
\qty{303.15}{K}$~: \qty{349}{m/s}. Pour \qty{9}{s}~: $9 \times 331 \approx
\qty{2980}{m}$ et $9 \times 349 \approx \qty{3140}{m}$, une différence
d'environ \qty{160}{m}, environ \qty{5}{\percent}. La règle donne
\qty{3}{km}~: c'est assez pour savoir à quelle distance est l'\omterm{def:g10:storms-and-cyclones:thunderstorm}{orage}, et elle
ne sert qu'à cela.
\end{solution}

\begin{solution}{exo:g10:storms-and-cyclones:13}
$\Delta p = 1010 - 940 = \qty{70}{hPa}$~: le bombement vaut environ \qty{0.7}{m}, la
part du vent $4.0 - 0.7 = \qty{3.3}{m}$, plus de quatre fois plus. À marée
haute, l'eau atteint $4.0 + 1.5 = \qty{5.5}{m}$ au-dessus du niveau moyen
de la mer.
\end{solution}

\begin{solution}{exo:g10:storms-and-cyclones:14}
\omterm{def:g10:storms-and-cyclones:tornado}{Tornade}~: un tour fait $2\pi \times 100 \approx \qty{628}{m}$, parcourus en
$628/90 \approx \qty{7}{s}$. Ouragan~: $2\pi \times 30\,000 \approx
\qty{188}{km}$, parcourus en $188\,000/60 \approx \qty{3140}{s}$, environ
\qty{52}{min}. La \omterm{def:g10:storms-and-cyclones:tornado}{tornade} tourne quelque 450 fois plus vite~: c'est un
tourbillon petit et très rapide, l'ouragan un tourbillon immense et plus
lent.
\end{solution}

\begin{solution}{exo:g10:storms-and-cyclones:15}
$24/3888 \approx \qty{0.6}{\percent}$~; $414/3888 \approx
\qty{10.6}{\percent}$~; $121/3888 \approx \qty{3.1}{\percent}$. Si la
latitude ne comptait pas, environ \qty{9}{\percent} des tempêtes, quelque
350, naîtraient à moins de \qty{5}{\degree} de l'équateur~; 24 seulement le
font, environ quinze fois moins. Près de l'équateur, la rotation de la Terre
est trop faible pour amorcer la rotation de la tempête.
\end{solution}

\begin{solution}{pb:g10:storms-and-cyclones:1}
\textbf{1.} $\pi \times (665\,000)^2 \approx \qty{1.39e12}{m^2}$.

\textbf{2.} $0.015 \times 1.39\times10^{12} \approx \qty{2.08e10}{m^3}$.

\textbf{3.} Un kilomètre cube vaut \qty{e9}{m^3}~: environ \qty{21}{km^3}.

\textbf{4.} $2.08\times10^{10} \times 1000 \approx \qty{2.08e13}{kg}$.

\textbf{5.} $2.08\times10^{13} \times 2.5\times10^6 \approx \qty{5.2e19}{J}$.

\textbf{6.} $5.2\times10^{19}/86\,400 \approx \qty{6.0e14}{W}$.

\textbf{7.} $6.0\times10^{14}/1.22\times10^{17} \approx 0.005$, environ
\qty{0.5}{\percent}.

\textbf{8.} $240 \times 1.39\times10^{12} \approx \qty{3.3e14}{W}$. La tempête
libère presque deux fois plus de chaleur que la lumière solaire absorbée par
toute son aire.

\textbf{9.} La tempête ne fonctionne pas avec la lumière solaire qui tombe
sur elle~: elle libère de la chaleur stockée sous forme de vapeur d'eau,
évaporée sur une aire de mer chaude plus vaste et aspirée par les vents. La
tempête concentre de l'énergie recueillie sur une région bien plus grande.

\textbf{10.} La même masse que la pluie, environ \qty{2.1e13}{kg} par jour
(la vapeur apportée doit se renouveler aussi vite qu'elle tombe en pluie).

\textbf{11.} $6.0\times10^{14}/82 \approx \qty{7.3e12}{m^2}$, environ cinq fois
l'aire de la tempête~: la tempête doit puiser de la vapeur sur une aire bien
plus vaste, et ses vents font évaporer la mer sous elle bien plus vite que
la moyenne.

\textbf{12.} \qty{24.1}{g/kg} à \qty{28}{\degreeCelsius} contre
\qty{14.7}{g/kg} à \qty{20}{\degreeCelsius}~: $24.1/14.7 \approx 1.64$,
environ \qty{64}{\percent} de plus.

\textbf{13.} L'eau chaude s'évapore vite et l'air chaud contient bien plus
de vapeur~: l'air qui afflue apporte bien plus de \omterm{def:g10:earth-energy-budget:latent-heat}{chaleur latente} à libérer
et, saturé, il monte librement. Au-dessus de mers plus fraîches, l'apport
est trop faible pour faire tourner le moteur.

\textbf{14.} Sur la terre ferme, il n'y a pas de mer à évaporer~: l'apport
de \omterm{def:g10:earth-energy-budget:latent-heat}{chaleur latente} cesse, et le sol rugueux freine les vents.

\textbf{15.} $1.5\times10^{12} \times 86\,400 \approx \qty{1.3e17}{J}$.

\textbf{16.} $96 \times 1852/3600 \approx \qty{49.4}{m/s}$.

\textbf{17.} $\frac12 \times 1.16 \times 49.4^2 \approx \qty{1.4e3}{J}$.

\textbf{18.} $5.2\times10^{19}/1.3\times10^{17} = 400$~: la \omterm{def:g10:earth-energy-budget:latent-heat}{chaleur latente}
vaut quelque 400 fois l'énergie fournie aux vents.

\textbf{19.} L'essentiel réchauffe l'air d'altitude de la tempête, qui
s'étale au sommet et cède la chaleur à l'espace sous forme de rayonnement
infrarouge~; les vents eux-mêmes perdent leur énergie par frottement contre
la mer et le sol.

\textbf{20.} $1/400 = \qty{0.25}{\percent}$ (le script donne un rapport de
402)~: \textbf{environ une part sur 400 de la chaleur libérée par la
condensation devient du vent}.
\end{solution}
